\lecture{4}{16. Februar 2026}{Equilibrium Equations \& Principle of Angular Momentum}

\begin{problem}{2.2}
  The state of stress at a point $P$ in a structure is given by
  \begin{align*}
    \sigma_{11} &= \num{20000}, & \sigma_{22} &= \num{-15000} , & \sigma_{33} &= \num{3000} \\
    \sigma_{12} &= \num{2000} , & \sigma_{23} &= \num{2000}, & \sigma_{31} &= \num{1000}
  .\end{align*}

  \begin{enumerate}[label=(\alph*)]
    \item Compute the scalar components $T_1, T_2, $ and $T_3$ of the traction $\Vec{T}$ on the plane passing through $P$ whose outward normal vector $\Vec{n}$ makes equal angles with the coordinate axes $x_1$, $x_2$, $x_3$.
    \item Compute the normal and tangential components of stress on this plane.
    \item[(d)] Determine the principal stresses.
    \item[(e)] Determine the direction of maximum principal stress.
  \end{enumerate}
\end{problem}

\begin{problem}{2.3}
  The state of stress at a point is given by the components of stress tensor $\sigma_{ij}$.
  A plane is defined by the direction cosines of the normal $\left( 1 / 2, 1 / 2, 1 / \sqrt{2} \right)$.
  State the general conditions for which the traction on the plane has the same direction as the $x_2$-axis and a magnitude of \num{1,0}.
\end{problem}

\begin{problem}{2.4}
  Determine the body forces for which the following stress field describes a state of equilibrium
  \begin{align*}
    \sigma_{x x} &= - 2x^2 - 3y^2 - 5z, & \sigma_{xy} &= x + 4xy -6, \\
    \sigma_{yy} &= -2y^2 + 7, & \sigma_{xz} &= -3x + 2y + 1. \\
    \sigma_{zz} &= 4x + y + 3z - 5, & \sigma_{yz} &= 0
  .\end{align*}
\end{problem}

\begin{problem}{2.5}
  Determine whether the following stress field is admissible in an elastic body when body forces are negligible.
  \begin{equation*}
    \Vec{\sigma} = \begin{bmatrix}
      yz+4 & z^2 + 2x & 5y+z\\
    \cdot & xz + 3y & 8x^3\\
    \cdot  & \cdot  & 2xyz\\
    \end{bmatrix}
  .\end{equation*}
  
\end{problem}

\begin{problem}{2.8}
  The stress tensor at a point $P(x,y,z)$ in a solid is
  \begin{equation*}
    \Vec{\sigma} = \begin{bmatrix}
      3 & 1 & 1\\
    1 & 0 & 2\\
    1 & 2 & 0\\
    \end{bmatrix}
  .\end{equation*}
  \begin{enumerate}[label=(\alph*)]
    \item Show that the principal stresses are \num{4} , \num{1} , and \num{-2} .
    \item Find the orientation of the principal planes
  \end{enumerate}
\end{problem}

\begin{problem}{2.9}
  The stress tensor at a point $P(x,y,z)$ in a solid is
  \begin{equation*}
    \Vec{\sigma} = \begin{bmatrix}
      \num{120}  & \num{40}  & \num{30} \\
    \num{40}  & \num{150}  & \num{20} \\
    \num{30}  & \num{20}  & \num{100} \\
    \end{bmatrix}
  .\end{equation*}
  
  \begin{enumerate}[label=(\alph*)]
    \item Determine the stresses with respect to a set of axes $x'$, $y'$, and $z'$ for which the matrix of direction cosines is
      \begin{equation*}
        \left[ \alpha_{ij} \right] = \begin{bmatrix}
          \frac{1}{2} & - \frac{\sqrt{3}}{2} & 0\\
        \frac{\sqrt{3}}{2} & \frac{1}{2} & 0\\
        0 & 0 & 1\\
        \end{bmatrix}
      .\end{equation*}
    \item Compute the principal stresses using both the $x-$, $y-$, and $z-$ and $x'-$, $y'-$, and $z'-$axes as the basis.
    \item Find the orientation of the principal planes with respect to the $x-$, $y-$, and $z-$axes.
  \end{enumerate}
\end{problem}
